\documentclass[11pt]{article}

\usepackage[utf8]{inputenc}
\usepackage{graphicx}
\usepackage[dvipsnames]{xcolor}
\usepackage{color}
\usepackage[margin=1in]{geometry}
\usepackage{hyperref}
\usepackage{tikz}
\usepackage{amsmath}
\usepackage{commath}
\usepackage{relsize}
\usepackage{centernot}
\usepackage{amssymb}
\usepackage{amsfonts}
\usepackage{graphicx}
\usepackage{textcomp}
\usepackage{gensymb}
\usepackage{enumitem}
\usepackage{siunitx}
\usepackage{fancyhdr}
\usepackage{floatrow}
\usepackage{wrapfig}
\usepackage[labelfont=bf]{caption}
\usepackage{mathtools}
\usepackage{pgfplots}
\usepackage{multicol}
\usepackage[misc]{ifsym}
\usepackage{tabularx}
\usepackage{empheq}
\usepackage{tkz-tab}
\usepackage{xpatch}
\usepackage{caption}
\usepackage{subcaption}
\usepackage{tabularray}
\usepackage{esint}
\usepackage{amsthm}
\usepackage{thmtools}
\RequirePackage[most]{tcolorbox}  

\swapnumbers

\makeatletter
\declaretheoremstyle[
  headfont= \sffamily\bfseries\color{cyan!50!black},
  headpunct={\sffamily\bfseries.},
  postheadspace=0.5em,
  notefont=\sffamily\bfseries,
  headformat=\NAME\NUMBER\let\thmt@space\@empty\NOTE,
  bodyfont=\mdseries\slshape,
  spaceabove=12pt,spacebelow=12pt,
]{thmstyle}
\makeatother

\theoremstyle{thmstyle}
\declaretheorem[name=Definición]{theorem}
\newtheorem{definition}[theorem]{Definición}

\definecolor{lightcyan}{rgb}{0.88, 1.0, 1.0}
\colorlet{mythmback}{lightcyan!40!white}

\tcolorboxenvironment{theorem}{
  colback=mythmback,coltitle=blue,colframe=mythmback,
  left=0pt,right=0pt,
  top=0pt,bottom=0pt,
  before skip=10pt, after skip=10pt,
}

\xpatchcmd{\tkzTabLine}{$0$}{$\bullet$}{}{}
\tikzset{t style/.style={style=solid}}
\pgfplotsset{compat=newest}

\newtheorem{manualtheoreminner}{}
\newenvironment{manualtheorem}[1]{%
  \renewcommand\themanualtheoreminner{#1}%
  \manualtheoreminner
}{\endmanualtheoreminner}

\newcommand{\pd}[2]{\frac{\partial #1}{ \partial #2}}
\newcommand{\td}[2]{\frac{\mathrm{d} #1}{ \mathrm{d} #2}}
\newcommand{\pdd}[2]{\frac{\partial^2 #1}{ \partial #2 ^2}}
\newcommand{\pddc}[3]{\frac{\partial^2 #1}{ \partial #2 \partial #3}}
\newcommand{\solution}[0]{\large{\textbf{Solución}}\normalsize}

\def\sca   #1{\mbox{\rm{#1}}{}}
\def\mat   #1{\mbox{\boldmath $\mathsf #1$}}
\def\vec   #1{\mbox{\boldmath $#1$}{}}
\def\ten   #1{\mbox{\boldmath $#1$}{}}
\def\ltr   #1{\mbox{\sf{#1}}}
\def\bltr  #1{\mbox{\sffamily{\bfseries{{#1}}}}}

\DeclareMathOperator{\dive}{div}
\DeclareMathOperator{\trace}{trace}
\DeclareMathOperator{\tr}{tr}
\DeclareMathOperator{\symm}{symm}
\DeclareMathOperator{\sk}{skew}
\DeclareMathOperator{\grad}{grad}
\DeclareMathOperator{\Grad}{Grad}
\DeclareMathOperator{\curl}{curl}
\DeclareMathOperator{\Curl}{Curl}
\DeclareMathOperator*{\argmin}{\arg\!\min}

\def\dx{\mbox{\(\,\mathrm{d}x\)}}
\def\ds{\mbox{\(\,\mathrm{d}s\)}}
\def\dS{\mbox{\(\,\mathrm{d}S\)}}
\def\dV{\mbox{\(\,\mathrm{d}V\)}}
\def\dv{\mbox{\(\,\mathrm{d}v\)}}
\def\R{\mathbb R}
\def\C{\mathbb C}
\def\N{\mathbb N}
\def\Q{\mathbb Q}
\def\Z{\mathbb Z}

\usepackage{mdframed}
\mdfdefinestyle{MyFrame}{%
    linecolor=black,
    linewidth=1pt,
    outerlinewidth=1pt,
    roundcorner=20pt,
    innertopmargin=\baselineskip,
    innerbottommargin=\baselineskip,
    innerrightmargin=20pt,
    innerleftmargin=20pt,
    splittopskip=2\baselineskip,
    backgroundcolor=gray!50!white}
    
\setlength\textwidth{7in}
\setlength\parindent{0pt}
\oddsidemargin=-0.5cm
\pagestyle{fancy}
\newcommand*{\vertbar}{\rule[-1ex]{0.5pt}{2.5ex}}
\setlength{\headheight}{13.59999pt}

\renewcommand{\headrulewidth}{0.1pt}
\renewcommand{\footrulewidth}{0.1pt}
\renewcommand{\figurename}{Figura}
\renewcommand\refname{Referencias}
\renewcommand{\thesection}{\arabic{section}.}
\renewcommand{\thesubsection}{\thesection\arabic{subsection}.}
\renewcommand{\labelitemi}{{\raisebox{.3\height}{\scalebox{0.6}{$\blacklozenge$}}}}
\renewcommand*\contentsname{Índice}

\newcommand{\p}{\vspace{10pt}}

\AtBeginDocument{\vspace*{-3.2\baselineskip}}
\textheight=665pt

\lhead{\scshape Pontificia Universidad Católica de Chile}
\rhead{\scshape IMC UC}
\cfoot{\thepage}

\usepackage{footnote}
\usepackage{etoolbox}
\newtoggle{indefintion}
\togglefalse{indefintion}
\pretocmd{\footnote}{\iftoggle{indefintion}{\stepcounter{footnote}}{\relax}}{}{}

\begin{document}
    \begin{flushleft}
        \small
        \vspace{1pt}
        \textbf{IMT3130} \hfill
        \textbf{13/03/2026}
    \end{flushleft}
    \begin{center}
        \LARGE\textbf{Ayudantía 1}\\
        \vspace{3pt}
        \normalsize\textbf{Cálculo diferencial tensorial}\\
        \normalsize Ayudante: Bastián Herrera \Letter{\hspace{1pt}: \texttt{biherrera@uc.cl}}
        \rule{\textwidth}{0.05mm}
    \end{center}

\section{Problemas}
\begin{enumerate}[label=\textbf{\arabic*.}, leftmargin=*]

    \item Sea $F:\R^n\to\R$ definida por
    \[
    F(\vec x)=\frac{1}{2}\vec x^\top \ten M \vec x+\vec b\cdot \vec x,
    \]
    donde $\ten M\in\R^{n\times n}$ y $\vec b\in\R^n$. Calcule la derivada de Fréchet de $F$ en un punto $\vec x\in\R^n$.

    \item Sea $\Omega\subset\R^d$ un dominio acotado, $\alpha>0$ y $f\in L^1(\Omega)$. Considere el funcional $J:L^\infty(\Omega)\to\R$ dado por
    \[
    J(u)=\int_\Omega\left(\frac{1}{2}u^2+\frac{\alpha}{4}u^4-fu\right)\dx.
    \]
    \begin{enumerate}[label=(\alph*)]
        \item Calcule la derivada de Gâteaux de $J$ en la dirección $v\in L^\infty(\Omega)$.
        \item Verifique que $J$ es Fréchet diferenciable y determine explícitamente su derivada.
    \end{enumerate}

    \item Sean $\ten A,\ten H,\ten B\in\R^{n\times n}$, con $\ten B$ fijo.
    \begin{enumerate}[label=(\alph*)]
        \item Defina $\ten C(\ten A)=\ten A^\top \ten A$. Calcule $d\ten C(\ten A)[\ten H]$.
        \item Defina $\psi(\ten A)=\ten A:\ten B=\tr(\ten A^\top \ten B)$. Calcule $d\psi(\ten A)[\ten H]$ y la derivada tensorial $\pd{\psi}{\ten A}$.
    \end{enumerate}

    \item Para $\ten A\in\R^{n\times n}$, considere las invariantes
    \[
    I_1(\ten A)=\tr(\ten A), \qquad
    I_2(\ten A)=\frac{1}{2}\left[(\tr \ten A)^2-\tr(\ten A^2)\right], \qquad
    I_3(\ten A)=\det(\ten A).
    \]
    Calcule las derivadas tensoriales
    \[
    \pd{I_1}{\ten A},\qquad \pd{I_2}{\ten A},\qquad \pd{I_3}{\ten A},
    \]
    suponiendo en el último caso que $\ten A$ es invertible.

    \item Sea $W:\{\ten F\in\R^{n\times n}:\det\ten F>0\}\to\R$ la función dada por
    \[
    W(\ten F)=\frac{\mu}{2}\big(\ten F:\ten F-3\big)-\mu\ln(\det\ten F)+\frac{\lambda}{2}\big(\ln(\det\ten F)\big)^2,
    \]
    donde $\mu,\lambda>0$.
    \begin{enumerate}[label=(\alph*)]
        \item Calcule la derivada tensorial $\pd{W}{\ten F}$.
        \item Evalúe el resultado en $\ten F=\ten I$.
    \end{enumerate}

    \item En los siguientes ejemplos, analice la existencia de derivadas direccionales o de Gâteaux en el origen y determine si la función es Fréchet diferenciable en dicho punto.
    \begin{enumerate}[label=(\alph*)]
        \item Considere $f:\R\to\R$ dada por
        \[
        f(x)=|x|.
        \]
        Estudie la diferenciabilidad direccional y la diferenciabilidad de Fréchet en $x=0$.
        \item Considere $f:\R^2\to\R$ definida por
        \[
        f(x,y)=
        \begin{cases}
        \dfrac{2xy}{\sqrt{x^2+y^2}}, & (x,y)\neq(0,0),\\[1ex]
        0, & (x,y)=(0,0).
        \end{cases}
        \]
        Estudie la continuidad en $(0,0)$, calcule las derivadas de Gâteaux en ese punto y determine si $f$ es Fréchet diferenciable en $(0,0)$.
    \end{enumerate}

\end{enumerate}

\end{document}